\section{Citação}

\paragraph{}Em um trabalho científico devemos ter sempre a preocupação de fazer referências precisas às idéias, frases ou conclusões de outros autores, isto é, citar a fonte (livro, revista e todo tipo de material produzido gráfica ou eletronicamente) de onde são extraídos esses dados. As citações fundamentam e melhoram a qualidade científica do trabalho, portanto, elas têm a função de oferecer ao leitor condições de comprovar a fonte das quais foram extraídas as idéias, frases ou conclusões, possibilitando-lhe ainda aprofundar o tema/assunto em discussão. Têm ainda como função, acrescentar indicações bibliográficas de reforço ao texto. Veja alguns exemplos:

\paragraph{}É neste cenário que "[...] o desafiador cenário globalizado cumpre um papel essencial na formulação das formas de ação." \cite{Papadimitriou04}.

\paragraph{}Segundo Flávio Mello \cite{BrainVoyager06}, "[...] é claro que o início da atividade geral de formação de atitudes prepara-nos para enfrentar situações atípicas decorrentes do sistema de formação de quadros que corresponde às necessidades [...]".

\paragraph{}De acordo com Flávio Mello \cite{Heeger02}, a certificação de metodologias que nos auxiliam a lidar com o acompanhamento das preferências de consumo causa impacto indireto na reavaliação do orçamento setorial.

\paragraph{}Por outro lado, a mobilidade dos capitais internacionais garante a contribuição de um grupo importante na determinação dos índices pretendidos.  Percebemos, cada vez mais, que o novo modelo estrutural aqui preconizado deve passar por modificações independentemente das condições inegavelmente apropriadas \cite{Canguilhem80}.

\paragraph{}Além disto, a expressão latina \textbf{apud} que significa: citado por, conforme, segundo é utilizada quando se faz referência a uma fonte secundária. Suponha que você teve acesso ao conteúdo do texto de Fulado através do trabalho Beltrano:
Não obstante, a competitividade nas transações comerciais prepara-nos para enfrentar situações atípicas decorrentes do fluxo de informações. A determinação clara de objetivos garante a contribuição de um grupo importante na determinação dos relacionamentos verticais entre as hierarquias. O incentivo ao avanço tecnológico, assim como a mobilidade dos capitais internacionais desafia a capacidade de equalização de alternativas às soluções ortodoxas (Fulano \cite{Fulano03} apud Beltrano \cite{Beltrano99}).

\section{Figuras}

\paragraph{}Figuras (organogramas, fluxogramas, esquemas, desenhos, fotografias, gráficos, mapas, plantas e outros) constituem unidade autônoma e explicam, ou complementam visualmente o texto, portanto, devem ser inseridas o mais próximo possível do texto a que se referem. Sua identificação deverá aparecer na parte inferior precedida da palavra designativa (figura), seguida de seu número de ordem de ocorrência, do respectivo título e/ou legenda e da fonte, se necessário, tal como na Figura \ref{FigDel}.

\begin{figure}
\begin{center}
\parbox[htb]{13.0cm}
  {
  \begin{center}
  \includegraphics[scale=1.0]{logo_del.eps}
  \caption[\small{Logotipo do DEL. Fonte: DEL/Poli/UFRJ \cite{Meyer97}.}]{\label{FigDel} \small{Logotipo do DEL. Fonte: DEL/Poli/UFRJ \cite{Meyer97}.}}
  \end{center}
  }
\end{center}
\end{figure}

\section{Tabelas}
\paragraph{}As tabelas são elementos demonstrativos de síntese que apresentam informações tratadas estatisticamente constituindo uma unidade autônoma. Em sua apresentação deve ser observado: (1) o título deverá ser colocado na parte inferior, precedido da palavra Tabela e de seu número de ordem; (2) as fontes e eventuais notas aparecem em seu rodapé, após o fechamento, utilizando-se o tamanho 10; (3) Devem ser inseridas o mais próximo possível do trecho a que se referem, tal como a Tabela \ref{TabIntranet}.

\begin{table}[h]
	\begin{center}
	  \caption{Casos de ataques aos computadores da Intranet. Fonte: DEL/Poli/UFRJ \cite{Meyer97}.}
		\begin{tabular}{|c|c|c|}\hline
		  \textbf{Número I}P & \textbf{Ataques} & \textbf{Ataques bem sucedidos} \\ \hline \vspace{-1.0mm}
		  192.168.0.120 & 54 & 1 \\ \hline \vspace{-1.0mm}
		  192.168.0.123 & 36 & 2 \\ \hline \vspace{-1.0mm}
		  192.168.0.129 & 25 & 4 \\ \hline \vspace{-1.0mm}
		  192.168.0.130 & 16 & 0 \\ \hline \vspace{-1.0mm}
		  192.168.0.141 & 29 & 3 \\ \hline \vspace{-1.0mm}
		  \textbf{Total} & \textbf{160} & \textbf{10} \\ \hline
		\end{tabular}
	\end{center}
	\label{TabIntranet}
\end{table}

\section{Numeração de páginas}

\paragraph{}O aluno deve observar atentamente a numeração de páginas de seu projeto. A primeira parte deste modelo de projeto final, composta pela dedicatória, agradecimento, resumo, abstract, siglas, sumário, lista de figuras e lista de tabelas, é numerada seqüencialmente utilizando algarismos romanos minúsculos. As demais folhas, descritas na segunda parte deste modelo, são numeradas seqüencialmente utilizando algarismos arábicos.

\paragraph{}Contudo, exclusivamente para a segunda parte do modelo de projeto, é permitida uma numeração alternativa na qual o aluno poderá numerar as páginas por capítulo. Por exemplo, a primeira página deste Capítulo 2 - Informações Adicionais, poderia ser escrita como 2.1. Além disto, a página seguinte seria 2.2 e a presente página poderia ser escrita como 2.3. A página do Apêndice A - O que é um apêndice, poderia ser escrita como A.1, enquanto que a primeira página do apêndice B seria B.1. Neste caso alternativo específico, a Bibliografia na deverá conter numeração.

